\documentclass[10pt,a4paper]{amsart}
\usepackage[margin=22mm]{geometry}
\usepackage[no-math]{fontspec}\usepackage{amsmath,amssymb,mathtools,mathalfa,xfrac,xcolor}
\setmainfont{Latin Modern Roman}[SmallCapsFont=lmromancaps10-regular.otf]
\usepackage[colorlinks=true,linkcolor=blue,citecolor=teal,urlcolor=blue]{hyperref}
\sloppy\newcommand{\eg}{e.g.\ }\newcommand{\ie}{i.e.\ }\newcommand{\resp}{resp.\ }\newcommand{\cf}{cf.\ }\newcommand{\skpt}{}\let\Subsection\subsection\let\Subsubsection\subsubsection
\let\ts\textstyle
\let\epsilon\varepsilon
\newcommand{\quand}{\quad\text{and}\quad}
\newtheorem{theorem}{Theorem}[section]
\newtheorem{proposition}[theorem]{Proposition}
\newtheorem{lemma}[theorem]{Lemma}
\newtheorem{corollary}[theorem]{Corollary}

\newcommand{\R}{\mathbb{R}}
\newcommand{\Z}{\mathbb{Z}}
\newcommand{\N}{\mathbb{N}}
\newcommand{\PP}{\mathbb{P}}

\newcommand{\CC}{\mathcal C}
\newcommand{\F}{\mathcal{F}}
\newcommand{\M}{\mathcal M}
\newcommand{\cS}{\mathcal{S}}
\newcommand{\U}{\mathcal U}
\newcommand{\V}{\mathcal{V}}
\newcommand{\X}{\mathcal{X}}

\let\e\varepsilon
\let\G\Gamma
\let\g\gamma
\let\pp\partial
\let\la\lambda
\let\La\Lambda
\let\Om\Omega
\let\s\sigma

\let\x\times

\newcommand{\aaa}{\mathfrak a}

\DeclareMathOperator{\dist}{dist}
\DeclareMathOperator{\Gr}{Gr}
\DeclareMathOperator{\id}{Id}
\DeclareMathOperator{\nil}{Nil}
\DeclareMathOperator{\SL}{SL}
\begin{document}
\title{Falconer--Lyapunov dimension variational principle for Anosov representations}
\author{François Ledrappier\and Pablo Lessa}
\date{}\maketitle
\section{Introduction}

\subsection{Main results}

Let \((\Gamma,\cS)\) be a countable group with a symmetric finite set of generators \(S\).

For \(d \ge 2\), consider \(P\) a subset of \(\lbrace 1,\ldots, d-1\rbrace\).
A representation \(\rho: \Gamma \to \SL(d,\R)\)
is called a {\it{\(P\)-Anosov}} representation if there exists \(c > 0\) such that
\[
\frac{s_{p}(\rho(\gamma))}{s_{p+1}(\rho(\gamma))} > c\exp(c|\gamma|),
\]
for all \(\gamma \in \Gamma\) and \(p \in P\). Here \(|\gamma|\) denotes word length in \(\Gamma\) with respect to \(S\) and, for~\(1\le i\le d, \, s_i(\g)\) denotes the \(i\)-th largest singular value of \(\g \in \SL(d,\R)\) with respect to the usual inner product on \(\R^d\). A representation is called {\it{Borel-Anosov}} if it is \(\lbrace 1,\ldots, d-1\rbrace\)-Anosov.

Anosov representations have been introduced by F.\,Labourie (\cite{Lab06}) and have been the subject of many studies (\cf in particular \cite{GGKW17}, \cite {KLP16}, \cite{KLP17}, \cite{KLP18a}, \cite{BPS19} and \cite{canary} for a recent survey). It is widely accepted that non-trivial Anosov representations present the right analog of convex cocompact representations in higher rank. Many properties are suggested by this analogy. In another direction, we~present in this note a phenomenon which is purely higher rank: there is a dimension gap for minimal invariant subsets for the action of \(\rho (\G) \) on the spaces of flags. The proof uses a new variational principle for the action of \(\rho (\G) \) on the spaces of partial flags (see below Theorem \ref{theoremb}).

Let \(Q\) be a non-empty subset of \(\lbrace 1,\ldots, d-1 \rbrace\). We~consider the space \(\F_Q\) of partial flags with signature \(Q\), endowed with a rotationally invariant Riemannian metric. Let~us recall the definition of the Minkowski dimension of a set \(\La \subset \F_Q\): for \(\e >0\), let~\(N(\La, \e) \) be the covering number of the set \(\La \) by balls of radius \(\e\) in the metric of~\(\F_Q\). The {\it {Minkowski (or Box) dimension }} \(\dim_B(\La) \) is defined by
\[
 \dim_B (\La) \; : =\; \limsup_{\e \to 0} \frac{\log N(\La, \e)}{\log (1/\e)}.
\]
Recall that the Hausdorff dimension of a metric set is bounded above by its Box dimension. Our result is
\begin{theorem}[Dimension gap for minimal sets of Anosov representations]\label{main}
With the above notations, let \(\La\) be a minimal invariant set for the action of \(\rho (\G) \) on \(\F_Q\). Assume that the representation \(\rho\) is \(P\)-Anosov for some \(P \subset \lbrace 1,\ldots, d-1 \rbrace\) and such that \(\rho (\G)\) is Zariski dense in \(\SL(d,\R)\). Then, with \(M := \# (Q \cap P)\),
\[
\dim_B(\Lambda) \le \dim(\F_Q) - \frac{M - 1}{2}.
\]
\end{theorem}

In the case \(d = 2\) the bound given by the above theorem is trivial. And, indeed, uniform discrete subgroups of \(\SL(2,\R)\) have full limit set. A dimension gap for the limit set of classical Schottky groups was proved in \cite{Doy88}. We~note, however, that there exist uniform discrete (and thus convex co-compact) subgroups of isometries of
rank one symmetric spaces and thus there is no dimension gap for convex cocompact representations in rank one. On the other hand, there are many conditions that ensure that the limit set is a Lipschitz submanifold of the flag space (see the discussions in \cite{PSW19}). Our result is easy in those cases.

If \(\rho\) is \(P\)-Anosov then it is also \(P \cup P^\ast\)-Anosov for \(P^\ast = \lbrace d-p: p \in P\rbrace\). Theorem~\ref{main} gives a non-trivial upper bound for the dimension of any minimal invariant set in \(\F_{Q}\) for all \(Q\) with \(\#(Q\cap (P \cup P^\ast)) \ge 2\). In particular, as soon as the representation is \(k\)-Anosov for some \(k \not = d/2\) and Zariski dense, then the codimension of the limit set in the full flag space is at least 1/2.

There are a few simple reductions which allow us to state our core result. Firstly, we~observe that the bundle \(\pi_{Q, Q\cap P} : \F_Q \to \F_{Q\cap P} \) is smooth, so that for any closed \(\La \subset \F_{Q\cap P}, \) we have
\[
 \dim_B ((\pi_{Q, Q\cap P})^{-1} (\La)) \; = \; \dim_B (\La) + \dim \F_Q - \dim \F_{Q\cap P}.
\]
So, it suffices to show Theorem \ref{main} for the case \(Q= P\). In that case, the unique minimal invariant subset of \(\F_{P}\) is the limit set \(\La_\rho\).

By considering \(\G /{\textrm {Ker}} \rho\), we may assume that the representation \(\rho\) is faithful. Furthermore, we~may assume that \(\Gamma\) is non-elementary hyperbolic since, on the one hand, an Anosov subgroup of \(\SL (d,\R) \) is hyperbolic (\cite{KLP18}, \cite[Th.\,3.2]{BPS19}) and, on~the other hand, the limit set of an Anosov representation of an elementary (\ie virtually cyclic) hyperbolic group is finite. Taking all this into account, we~only have to prove
\begin{theorem}[Dimension gap for limit sets of Anosov representations]\label{maintheorem}
Let \(\G\) be a non-elementary, hyperbolic, finitely generated group, \(P \subset \lbrace 1,\ldots, d -1\rbrace\) with \(M:= \#P \ge 2 \), \(\rho \) a faithful \(P\)-Anosov representation of \(\G\) in \(\SL (d,\R) \) with \(\rho(\G)\) Zariski dense in \(\SL (d,\R) \) and \(\La_\rho \subset \F_{P} \) the limit set of the representation \(\rho\). Then,
\[
\dim_B(\Lambda_\rho) \le \dim(\F_{P}) - \frac{M - 1}{2}.
\]
\end{theorem}

It seems natural to conjecture that a dimension gap will exist between the limit set and the corresponding space of flags, for any Zariski dense representation into a reductive algebraic Lie group of rank \(r \ge 2\) which is Anosov with respect to at least two simple roots.

Let \(g\in \SL(d,\R)\). The numbers \(\la_i (g) := \lim_n \frac{1}{n }\log s_i(g^n) \) are the logarithms of the moduli of the eigenvalues of \(g\). Apply the Oseledets theorem in the trivial case of one matrix
\(g \in \SL (d,\R)\). The Oseledets decomposition is obtained from the Jordan decomposition by grouping together the spaces corresponding to the eigenvalues with the same modulus. These moduli are the \(\exp \la_i (g)\).

In our setting, if~\(\g \in \G\) (but a finite number), for all \(p \in P, \; \la_p (\rho(\g)) -\la_{p+1}(\rho(\g)) \geq c |\g| \). We~associate to every such \(\g \in \G \) an element \(\eta (\rho (\g))\) of the space \(X_P\) of decompositions of \(\R^d = E_1 \oplus \ldots \oplus E_{M+1} \), with \(\dim E_k = p_k- p_{k-1}\), namely
\[
 \eta (\rho (\g)) \; = \; E_1(\rho (\g)) \oplus \ldots \oplus E_{M+1} (\rho (\g)), 
\]
with the property that for all \(k = 1, \ldots, M+1\), the space \(E_{k} (\rho (\g)) \) is generated by the Jordan spaces of \(\rho(\g)\) such that the modulus \(\exp \la\) of the eigenvalue satisfies\footnote{With the convention that \(p_0=0, \la_{p_0} =+\infty\) and \(p_{M+1}:= d\).}
\[
\exp \la_{p_{k-1}}(\rho (\g)) >\exp \la \ge \exp \la_{p_{k}}(\rho (\g)).
\]
Observe that the space \(X_P\) can also be seen as the open \(\SL (d,\R) \)-invariant set of pair of flag spaces in general position in \(\F_P \x \F_{P^\ast} \).

\begin{corollary}\label{eigenvectors} Let \(\G\) be a hyperbolic finitely generated group, \hbox{\(P \!\subset\! \lbrace 1,\ldots, d-1 \rbrace\)} with \(M:= \# P \).
With the above notations, if~\(\rho \) is a \(P\)-Anosov, Zariski dense representation of \(\G\) in \(\SL(d,\R) \), set \(\Om_\rho \subset X_P \) the closure of the set \(\{ \eta (\rho (\g)); \g \in \G\}\). Then,
\[
\dim_B(\Om_\rho) \le \dim X_P -M +1.
\]
\end{corollary}

Indeed, Corollary \ref{eigenvectors} is trivial for \(d= 2\) and for \(M=1\). It is trivial as well if \(\G\) is elementary. For \(d \ge 3\) and \(M\ge 2\), we show in Section \href{https://jep.centre-mersenne.org/item/10.5802/jep.285.pdf#page=13}{2.3} why Corollary \ref{eigenvectors} follows from Theorem \ref{maintheorem}.

The first non-trivial case for our results is when \(d=3\) and \(P= \{1,2\}\). So,~let~\(\rho \) be a Borel-Anosov representation of the finitely generated non-elementary group \((\G, S)\) in \(\SL(3,\R) \) and we can consider the limit set \(\Lambda_\rho \) of the action of \(\rho(\G) \) on the space~\(\F\) of complete flags in \(\R^3\). Besides, for all \(\g \in \G\), the matrix \(\rho(\g)\) admits three distinct eigenspaces written \((E_1 (\rho(\g)), E_2 (\rho(\g)), E_3 (\rho(\g)))\) in the order of the absolute values of the eigenvalues. Let \(\Om_\rho \) denote the closure in \(\PP^2\x\PP^2\x\PP^2\) of the \(\{(E_1 (\rho(\g)), E_2 (\rho(\g)), E_3 (\rho(\g))), \g \in \G \} \). We~obtain, from Theorem \ref{maintheorem} and Corollary~\ref{eigenvectors}:

\skpt
\begin{corollary}\label{2.5, 5}\mbox{}\vspace*{-\baselineskip}
\[
\dim_B(\Lambda_\rho) \le \frac{5}{2} < 3 = \dim(\F), \; \dim_B(\Om_\rho) \le 5 < 6 = \dim (\PP^2\x\PP^2\x\PP^2).
\]
\end{corollary}
Let \(\G\) be a surface group. For the Hitchin component, the dimension of the projective limit set is 1 (see \cite{Lab06}) and our result is not new. It is is new and unexpected, as far as the authors are aware, for the Zariski dense representations that lie in the same connected component as the Teichmüller times Identity representations, \ie in the Barbot component (\cf \cite{barbot}).

\Subsection{Strategy of the proof of Theorem \ref{maintheorem}}

\subsubsection{Random walk entropy}

Let \(\M\) be the set of probability measures \(\mu\) on \(\Gamma\) with countable support and with finite first moment, meaning
\[
\sum_{\gamma \in \Gamma}\mu(\gamma) |\gamma| < +\infty,
\]
such that the semi-group \(\Gamma_\mu\) generated by the support of \(\mu\) is non-elementary (\ie contains two independent loxodromic elements).
By \cite{KV83}, if~\(\mu \in \M\) then \(\mu\) has finite Shannon entropy, meaning
\[
H(\mu) = -\sum_{\gamma \in \Gamma}\mu(\gamma)\log(\mu(\gamma)) < +\infty.
\]
The random walk entropy \(h_\mu\) of \(\mu \in \M\) is defined by
\begin{equation}\label{rwentropy} h_\mu = \lim_{n \to +\infty}\frac{1}{n}H(\mu^{\ast n}),\end{equation}
where \(\mu^{\ast n}\) is the \(n\)-fold convolution of \(\mu\) with itself.

\subsubsection{Lyapunov exponents}\label{lyapunovexp}
Denote by \(\aaa^+ \) the cone
\[
\ts \aaa^+ := \lbrace a \in \R^{d}: a_1 \ge \dots \ge a_d, \sum a_i = 0\rbrace.
\]
For \(g \in \SL(d,\R)\) one has \(\log S(g) \in \aaa^+\), where
\[
\log S(g): = (\log s_1(g),\ldots, \log s_d(g)). 
\]
The Lyapunov exponents \(\lambda(\rho_*\mu) = \{\lambda_i(\rho_\ast\mu), 1 \le i \le d \} \in \aaa^+\) of the random walk \((\G, \mu)\) in the representation \(\rho\) are defined by
\begin{equation}\label{exponents}
\lambda(\rho_*\mu):= \lim_{n \to +\infty} \frac{1}{n}\sum_{\gamma \in \Gamma}\mu^{\ast n}(\gamma)\log\bigl(S(\rho(\gamma))\bigr).\end{equation}
The limits exist by the sub-additivity of the sequences
\[
\ts \{\sum_{j\le i}\sum_{\gamma \in \Gamma}\mu^{\ast n}(\gamma)\log\bigl(s_j(\rho(\gamma))\bigr) \}_{n \in \N}.
\]

\subsubsection{Lyapunov and Falconer dimension}\label{Lyadimension}

We say a pair \((i,j)\), \(0\!<\!i\!<\!j\!\le\! d\), is~separated by \(p \in P\) if \(i \le p < j\). For each \(p \in P\) let \(S_p\) be the set of pairs separated by \(p\), and define \(S(P) = \bigcup_{p \in P}S_p\). On the cone \(\aaa^+\) we define the roots
\[
\alpha_{i,j}(a) = a_i - a_j,
\]
for \(i < j\) and notice that they are non-negative on \(\aaa^+\). Lyapunov and Falconer dimensions are special values of Lyapunov and Falconer functionals on \(\aaa^+\).

We define, for \(h \ge 0\) the Lyapunov functional on \(\aaa^+\) by
\begin{align*}L^P_h(a) &= \text{maximum of } \sum_{(i,j) \in S(P)}r_{i,j}
\\ &\hphantom{{}={}}\text{subject to }0 \le r_{i,j} \le 1
\quand \sum_{(i,j) \in S(P)} r_{i,j}\alpha_{i,j}(a) \le h.
\end{align*}
For \(r \ge 0\) we also define the Falconer functional on \(\aaa^+\) by
\begin{align*}F^P_r(a) &= \text{minimum of } \sum_{(i,j) \in S(P)}r_{i,j}\alpha_{i,j}(a)
\\ &\hphantom{{}={}}\text{subject to }0 \le r_{i,j} \le 1
\quand \sum_{(i,j) \in S(P)} r_{i,j} \ge r.
\end{align*}

\begin{lemma}[Duality between Falconer and Lyapunov functionals]\label{lyapunovvsfalconer}
Given \(a \in \aaa^+\), for all \(r,h \ge 0\) one has \(F_r^P(a) \le h\) if and only if \(L_h^P(a) \ge r\).

Furthermore, let \(d(a)\) be the number of pairs \((i,j) \) in \(S(P)\) with \(\alpha_{i,j}(a) = 0\).
Then, \(F^P_r(a) = 0 \) for \(r \in [0,d(a)]\) and \(r \mapsto F^P_r(a)\) is an increasing homeomorphism from \([d(a),D]\) to \([0,F_{D}^P(a)]\) where \(D = \# S(P)=\dim (\F_P)\), and \(h \mapsto L_h^P(a)\) is its inverse.
\end{lemma}
\begin{proof}
To establish the first claim observe that directly from the definitions, both conditions \(F_r^P(a) \le h\) and \(L_h^P(a) \ge r\) are equivalent to there being some choice of \(r_{i,j} \in [0,1]\) for \((i,j) \in S(P)\) such that
\[
\sum_{(i,j) \in S(P)} r_{i,j}\alpha_{i,j}(a) \le h\quand\sum_{(i,j) \in S(P)} r_{i,j} \ge r.
\]
Moreover, in computing \(F^P_r(a)\), we can always take \(r_{i,j} = 1 \) for the pairs \((i,j) \) in \(S(P)\) with \(\alpha_{i,j}(a) = 0\). It follows that \(F^P_r(a) = 0 \) for \(r \in [0,d(a)]\). Then, there exists \(\e >0\) with \(\alpha_{i,j}(a) > \e\) for all the other pairs \((i,j)\). It follows immediately that \(F^P_{r+s}(a) >F^P_r(a) + \e s\) for all \(d(a) \le r < r+s \le D\). Hence \(r \mapsto F^P_r(a)\) is increasing on \([d(a), D]\). Since \(F_{d(a)}^P(a) = 0\), \(r \mapsto F^P_r(a)\) it is an increasing homeomorphism as claimed.

Similarly, picking \(\e > 0\) small enough so \(\alpha_{i,j}(a) < \e^{-1}\) for all \((i,j) \in S(P)\), it~follows that \(L^P_{h + s}(a) > L^P_h(a) + \e s\) for all \(0 \le h < h+s < F_D^P(a)\). Hence, \(h \mapsto L_h^P(a)\) is an increasing homeomorphism as well.

Let \(f:I \to J\) and \(g:J \to I\) be the two homeomorphisms under consideration where \(I = [d(a),D]\) and \(J = [0,F_D^P(a)]\). By our first claim we have \(g(f(r)) \ge r\) for all \(r \in I\) and \(f(g(h)) \le h\) for all \(h \in J\). However, substituting \(r = f^{-1}(h)\), the first inequality implies \(g(h) \ge f^{-1}(h)\) for all \(h \in J\), from which we obtain \(f(g(h)) \ge h\) for all \(h \in J\). Hence, we~have \(f(g(h)) = h\) for all \(h \in J\) as required.
\end{proof}

Following \cite{DO80} and \cite{KY79}, we~define the Lyapunov dimension of \(\rho_*\mu\) on \(\F_P\) by
\begin{equation}\label{dimlyap:definition}\dim_{LY}^P(\rho_*\mu) \; := \; L_{h_\mu}^P (\lambda(\rho_*\mu)). \end{equation}
In the spirit of \cite{Fal88}, we~define the Falconer dimension \(\dim_F^P(\rho)\) of \(\rho\) relative to~\(\F_P\) as the critical parameter \(r \ge 0\) for convergence of the series
\begin{equation}\label{falconerseries} \Phi_\rho^P (r) \; := \; \sum_{\gamma \in \Gamma} \exp \bigl(-(F_r^P\circ\log S(\rho(\gamma)))\bigr).\end{equation}

\subsubsection{Proof of Theorem \ref{maintheorem}}

Theorem \ref{maintheorem} is a direct consequence of the three following results, all under the assumptions of Theorem \ref{maintheorem} and using the above notations:

\begin{theorem}[Upper bound on Lyapunov dimension]\label{theorema}
For all \(\mu \in \M\), one has
\[
\dim_{LY}^P(\rho_*\mu) \le \dim(\F_P) - \frac{M - 1}{2}.
\]
\end{theorem}

\begin{theorem}[Variational principle]\label{theoremb}
One has \begin{equation} \label{Lyapunov} \dim_F^P(\rho) = \sup_{\mu \in \M}\dim_{LY}^P(\rho_*\mu).\end{equation}
\end{theorem}

\begin{theorem}[Inequality between Minkowski and Falconer dimension]\label{theoremc}
One has
\[
\dim_B(\Lambda_\rho) \le \dim_F^P(\rho).
\]
\end{theorem}

The key observation behind Theorem \ref{theorema} is that the entropy \(h_\mu\) is realized as the Furstenberg entropy of the action of \(\rho (\G)\) on the Grassmannian \(\Gr(p,\R^d)\) for any \(p \in P\). This follows from the \(P\)-Anosov property and the identification of the Poisson boundary and geometric boundary for \(\mu\)-random walks on \(\Gamma\) (\cite{Kai00}). The inequality in Theorem \ref{theorema} is then a consequence of the inequality between Furstenberg entropy and the sum of the relevant exponents (see \cite{LL23} and Section \href{https://jep.centre-mersenne.org/item/10.5802/jep.285.pdf#page=14}{3}).

Concerning Theorem \ref{theoremb}, the fact that
\begin{equation}\label{LyapleqFalc} \dim_{LY}^P(\rho_*\mu) \; \leq \; \dim_F^P (G) \end{equation} is classical. We~recall the proof in Section \ref{proofLyap}.

The converse is due to Yuxiang Jiao, Jialun Li, Wenyu Pan and Disheng Xu (\cite{JLPX23}) in the case when the representation is Borel-Anosov. In \cite{JLPX23}, given \(\e>0\), the authors construct a free semi-group in \(\G\) that is rich enough that the uniform probability measure \(\mu_\e \) on the generators satisfies \(\dim_{LY}(\rho_*\mu_\e) \ge \dim_F(\rho) -\e\). The complete dominated splitting is used to control the almost additivity of the singular values.
Here, we~consider a general subset \(P\) and we assume that the representation \(\rho(\G)\) is Zariski dense in \(\SL(d,\R)\). In the spirit of \cite{gouezel}, \cite{random} and \cite{JLPX23}, we~prove in the appendix a general proposition about free sub-semigroups in hyperbolic groups (see the appendix for the definition of a quasi-geodesic set; see~also~\cite{yang} for a construction of free sub-semi-groups with convexity properties in a more general context.)
\begin{proposition}[$={}$Proposition \ref{semigrouplemma}]
Let \((\Gamma,\cS)\) be a non-elementary hyperbolic group with a finite symmetric generator \(\cS\). Then there exists a semigroup \(S \subset \Gamma\), a~finite set \(F \subset \Gamma\) and mappings \(L,R:\Gamma \to F\) with the following properties:
\begin{enumerate}
\item \(S\) generates \(\Gamma\) as a group,
\item \(S\) is quasi-geodesic, and
\item \(L(\gamma)\gamma R(\gamma) \in S\) for all \(\gamma \in \Gamma\).
\end{enumerate}
\end{proposition}
\section{Random walk entropy}

Let \(\G\) be a non-elementary finitely generated word hyperbolic group, \(P \subset \lbrace 1,\ldots, d-1 \rbrace\) with \(M:= \#P \ge 2 \), \(\rho \) a \(P\)-Anosov representation of \(\G\) in \(\SL (d,\R) \). We~let \(\mathcal{M}\) denote the set of probability measures \(\mu\) on \(\Gamma\) with finite first moment
\[
\sum_{\gamma} |\gamma| \mu(\gamma) < +\infty,
\]
and such that the semi-group \(\Gamma_\mu\) generated by the support of \(\mu\) contains two independent loxodromic elements.

In this section, we~prove that the random walk entropy \(h_\mu \) given by \eqref{rwentropy} coincides with the Furstenberg entropy (see below) of the stationary measure on the Gromov boundary of \(\Gamma\) and its images on adequate flag spaces. We~also relate the Falconer dimension and the random walk entropy to the growth indicator defined by \cite{Qui02a}. To conclude the section, we~prove Corollary \ref{eigenvectors}.

\subsection{Furstenberg entropy}

Recall that a probability measure \(\nu\) on a compact space \(X\) on which \(\Gamma\) acts continuously is said to be \(\mu\)-stationary, where \(\mu\) is a probability measure on \(\Gamma\), if~and only if
\[
\nu = \sum_{\gamma \in \Gamma}\mu(\gamma)\gamma_*\nu.
\]

The Furstenberg entropy of a \(\mu\)-stationary measure is defined as
\[
\kappa(\mu,\nu) = \sum_{\gamma \in \Gamma} \mu(\gamma) \int_{X} \log\Bigl(\frac{d\gamma_*\nu}{d\nu}(\g x)\Bigr)d\nu(x),
\]
or \(+\infty\) if the Radon-Nikodym derivative in the integral does not exist for some \(\gamma\) in the support of \(\mu\).

In this section we show that for \(\mu \in \M\), the Furstenberg entropy of the natural \(\mu\)-stationary measures on the Gromov boundary \(\partial\Gamma\), the Grassmannian manifolds of \(p\)-dimensional subspaces of \(\R^d\) for \(p \in P\), and the space of partial flags \(\F_P\) all coincide with the random walk entropy \(h_\mu\).

This particular feature of Anosov representations is useful since the Furstenberg entropy is what occurs in the dimension formulas of \cite{LL23}.

\Subsubsection{Stationary measure on the Gromov boundary \(\partial \Gamma\)}

\begin{theorem}\label{furstenberg}
For each \(\mu \in \M\) there exists a unique \(\mu\)-stationary measure \(\nu_\mu\) on the Gromov boundary \(\partial \Gamma\) of \(\Gamma\), and its Furstenberg entropy is given by \(\kappa(\mu,\nu_\mu) = h_\mu\).
\end{theorem}
\begin{proof}
If \(\g_{-1},\ldots, \g_{-n},\ldots\) are i.i.d. random elements in \(\Gamma\) with common distribution \(\mu\). Then by \cite[Th.\,7.6]{Kai00} there exists a random limit point
\[
X = \lim_{n \to +\infty}\g_{-1}\cdots \g_{-n} \in \partial\Gamma,
\]
almost surely, and the distribution \(\nu_\mu\) of \(X\) is the unique \(\mu\)-stationary measure on \(\partial \Gamma\).

Furthermore, \((\partial \Gamma,\nu_\mu)\) is isomorphic to the Poisson boundary of \((\Gamma,\mu)\) by \cite[Th.\,7.7]{Kai00}.

The Furstenberg entropy \(\kappa(\mu,\nu_\mu)\) is equal to the difference between \(h_\mu\) and the entropy of the random walk \(\g_1, \g_1\g_2,\ldots\) conditioned on \(X\) (see \cite[Cor.\,2]{KV83}). However, because \((\partial \Gamma,\nu)\) is the Poisson boundary this conditional entropy is zero (\cite[Th.\,4.3 and 4.5]{Kai00}) and therefore \(h_\mu = \kappa(\mu,\nu_\mu)\) as claimed.
\end{proof}

\subsubsection{Boundary maps}

Let \(p \in P \) and \(\gamma \in \Gamma\) be such that \(s_p(\rho(\gamma)) > s_{p+1}(\rho(\gamma))\).
There exists a unique \(p\)-dimensional subspace \(\xi^p(\gamma) \subset \R^d\) on which \(p\)-dimensional volume is most contracted by \(\rho(\gamma)^{-1}\).

We let \(\xi(\gamma) = (\{0\} \!\subset\! \xi^{p_1}(\gamma) \!\subset\! \ldots \!\subset\! \xi^{p_{M}}(\gamma))\!\subset\! \R^d \in \F_P\), where \hbox{\(P = \lbrace p_1 < \cdots < p_{M} \rbrace\)}, whenever all \(\xi^p(\gamma)\) are well defined.
Since \(\rho\) is \(P\)-Anosov this is the case outside of a finite subset of \(\Gamma\).

We refer to \cite[Chap.\,7]{GdlH90} for basic properties of the Gromov compactification of \(\Gamma\). We~briefly recall that \(\Gamma\) with its word metric is a proper geodesic metric space.
The group \(\Gamma\) is word hyperbolic, so there exists \(\delta > 0\) such that for every geodesic triangle each side is contained in the \(\delta\)-neighborhood of the other two.
The Gromov boundary \(\partial \Gamma\) is the set of equivalence classes of word geodesic rays in \(\Gamma\), where two rays are equivalent if they are at bounded Hausdorff distance.
A basis of neighborhoods of a point \(x \in \partial \Gamma\) in the Gromov compactification \(\overline{\Gamma} = \Gamma \cup \partial \Gamma\) is defined by taking for each \(C > 0\) the set
\[
\ts N(x,C) = \lbrace x\rbrace \cup \lbrace y \in \overline{\Gamma}: \min_{n} |\alpha_n| > C\text{ for all geodesics }\alpha\text{ joining }x\text{ and }y\rbrace.
\]

\begin{proposition}[\cite{GGKW17}]\label{boundarymapstheorem}
The map \(\xi\) defined above extends continuously to the Gromov boundary \(\partial \Gamma\). Furthermore, \(\xi^p:\partial\Gamma \to \Gr(p,\R^d)\) is injective for each \(p \in P\).
\end{proposition}

See \cite[Prop.\,30.3]{canary} for a direct proof.

\subsubsection{Dynamical stationary measures}

Fix \(\mu \in \M\), let \(\g_{-1},\ldots, \g_{-n},..\). be i.i.d.\ random elements of \(\Gamma\) with common distribution \(\mu\), \(\chi_1 > \cdots > \chi_N\) be the distinct Lyapunov exponents of \(\mu\), and \(d_1,\ldots, d_N\) their multiplicities.
Let
\[
A = \lbrace d_1, d_1+d_2,\ldots, d_1+\cdots + d_{N-1} \rbrace.
\]
By Oseledets theorem (\cite{Ose68}), for each \(a \in A\), there is a random limit
\[
U_a = \lim_{n \to +\infty} \xi^a(\g_{-1}\cdots \g_{-n}),
\]
almost surely.

The collection \(U\) of \(U_a\) for \(a \in A\) is a random element in the space of flags of signature \(A\). Its distribution, and the projections of its distribution to coarser partial flag spaces are what were called dynamical stationary measures in \cite{LL23}. Recall from Section \ref{lyapunovexp} the definition of the Lyapunov exponents \(\lambda_i(\rho_*\mu)\). Since \(\lambda_p(\rho_*\mu) > \lambda_{p+1}(\rho_*\mu)\) for all \(p \in P\) we have that \(\xi_*\nu_\mu\) is the dynamical stationary measure on \(\F_P\), and \(\xi^p_*\nu_\mu\) is the dynamical stationary measure on the Grassmannian of \(p\)-dimensional subspaces of \(\R^d\) for each \(p \in P\).

\Subsubsection{Furstenberg entropy}

\begin{theorem}\label{entropytheorem}
For each \(\mu \in \M\) one has
\[
h_\mu = \kappa(\mu,\xi_*\nu_\mu) = \kappa(\mu,\xi^p_*\nu_\mu),
\]
for all \(p \in P, p\not = 0,d\).
\end{theorem}
\begin{proof}
This is immediate from the injectivity of the maps \(\xi^p\) given by Proposition~\ref{boundarymapstheorem}.
\end{proof}

\Subsection{Growth indicator function}
\subsubsection{Growth function and Falconer dimension}

Given \(a \in \aaa^+\) we define the growth indicator at \(a\) by
\begin{equation}\label{growth}
\psi_\rho(a) = \| a\| \inf_{\CC \ni a}
\limsup_{T \to \infty } \frac{1}{T}\log \# \{ \g \in \G, \log S(\rho(\g)) \in \CC, \|\log S(\rho(\g)) \| \leq T\}, \end{equation}
where the infimum is over all open subcones \(\CC \subset \aaa^+\) that contain \(a\). The definition does not depend on the choice of the norm \(\|{\cdot}\| \) on \(\aaa^+\) (\cite{Qui02a}).

\begin{theorem}[Quint, \cite{Qui02a}]\label{cone}
The growth indicator function \(\psi_\rho\) is concave and \(\lbrace \psi_\rho > 0\rbrace\) is the interior of the limit cone defined by
\[
\mathcal{L}_\rho = \lim_{T \to +\infty}\frac{1}{T}\log S(\rho(\Gamma)),
\]
where the limit is taken in the Hausdorff topology on closed sets.
\end{theorem}

Recall that we defined the Falconer dimension \(\dim_F^P(\rho)\) as the critical value of the series
\(r \mapsto \sum_{\gamma \in \Gamma}\exp\left(-F^P_r\circ \log S(\rho(\gamma)\right)\). We~have (see \cite[Lem.\,4.2]{sambarino}, \cite[Lem.\,III.1.3]{Qui02a}):

\begin{lemma}\label{falconer2}
Assume \(r\ge 0 \) is such that there is \(a \in \aaa^+ \) inside the interior of the limit cone \(\mathcal{L}_\Gamma \) with \(\psi_\rho (a) > F^P_r (a). \) Then,
\[
 r \, \le \, \dim_F^P (\rho).
\]
\end{lemma}

\subsubsection{Inequality between Lyapunov and Falconer dimensions\label{proofLyap}}

Recall that, given \(\mu \in\nobreak \mathcal{M}\) we defined the Lyapunov dimension by
\[
\dim_{LY}^P(\rho_\ast \mu) = L_{h_\mu}(\lambda(\rho_*\mu)).
\]
We will now prove relation \eqref{LyapleqFalc}, which is the following statement:
\[
\sup_{\mu \in \mathcal{M}}\dim_{LY}^P(\rho_\ast \mu) \le \dim_F^P(\rho).
\]
For this purpose we first need the following version of the fundamental inequality for random walks, which is essentially due to Guivarc'h.

\begin{lemma}[Fundamental inequality]
For all \(\mu \in \mathcal{M}\) one has
\[
h_\mu \le \psi_{\rho}(\lambda(\rho_*\mu)).
\]
\end{lemma}
\begin{proof}
Consider the sequence space \(\Omega = \Gamma^{\Z}\) with the probability measure \(m = \mu^\Z\) and let \(g_n:\Omega \to \Gamma\) be the projection onto the \(n\)-th coordinate.
By the subadditive ergodic theorem \cite{Ose68} one has
\[
\lambda(\rho_*\mu) = \lim_{n \to +\infty}\frac{1}{n}\log S(\rho(g_n(\omega)\cdots g_1(\omega))),
\]
for \(m\)-almost every \(\omega \in \Omega\).
From the Shannon-McMillan-Breiman theorem for random walk entropy \cite{KV83} we have
\[
h_\mu = -\lim_{n \to +\infty}\frac{1}{n}\log \mu^{\ast n}(g_n(\omega)\cdots g_1(\omega)),
\]
for \(m\)-a.e. \(\omega \in \Omega\), where \(\mu^{\ast n}\) is the \(n\)-fold self-convolution of \(\mu\) and also the distribution of the random product \(g_1 \cdots g_n\).

Fix \(\epsilon > 0\) and let \(A_n \subset \Gamma\) be the set of elements satisfying
\[
n\lambda(\rho_*\mu) - n\epsilon < \log S(\gamma) < n\lambda(\rho_*\mu) + n\epsilon,
\]
and \(B_n \subset \Gamma\) consist of those elements satisfying
\[
\mu^{\ast n}(\gamma) < \exp(-n(h_\mu - \epsilon)).
\]
From the almost sure equalities above we have
\[
\lim_{n \to +\infty}\mu^{\ast n}(A_n \cap B_n) = 1.
\]
In particular
\[
\# (A_n \cap B_n) > \mu(A_n \cap B_n)e^{n(h_\mu - \epsilon)} > \frac{1}{2}e^{n(h_\mu - \epsilon)}
\]
for all \(n\) large enough.

This implies that for all \(\epsilon >0\) and \(n\) large enough, there are at least \(\frac{1}{2}e^{n(h_\mu -\e)}\) elements \(\g\) of \(\G\) with \(\log(S(\rho(\gamma)))\) \(n \e\)-close to \(n\lambda (\rho_*\mu)\). Which immediately yields the desired lower bound on the value of the growth indicator \(\psi_{\rho}(\lambda(\rho_*\mu))\).
\end{proof}

We now complete the proof of the inequality between the Lyapunov and Faloner dimensions.

\begin{proof}[Proof of inequality \eqref{LyapleqFalc}]
Fix \(\mu \!\in\! \mathcal{M}\) and let \(r_0 \!=\! \dim_{LY}^P(\rho_\ast \mu) \!=\! L_{h_\mu}(\lambda(\rho_*\mu))\). If~\hbox{\(r_0 = 0\)} there is nothing to prove, so we suppose \(r_0 > 0\).

From the duality between Lyapunov and Falconer functionals (Lemma \ref{lyapunovvsfalconer}) we have
\(F_{r_0}(\lambda(\rho_*\mu)) = h_\mu\). Furthermore, \(h_\mu \le \psi_{\rho}(\lambda(\rho_{*}\mu)) \) by the fundamental inequality. Hence
\[
F_{r_0}^P(\la(\rho_*\mu)) \le \psi_\rho (\la(\rho_*\mu)).
\]

Because \(\rho\) is \(P\)-Anosov, \(\lambda_i(\rho_*(\mu)) - \lambda_j(\rho_*(\mu)) > 0\) for all \((i,j) \in S(P)\). In particular, for each \(0\le r < r_0\) we have \(F_{r}^P(\la(\rho_*\mu)) < F_{r_0}^P(\la(\rho_*\mu))\) whenever \(0 \le r < r_0\) and hence
\(r \le \dim_F^P(\rho) \) by Lemma \ref{falconer2}, which concludes the proof.
\end{proof}

\setcounter{section}{3}
\section{Proof of Theorem \ref{theoremb}}\label{prooftheoremb}
The inequality \eqref{LyapleqFalc},
\(\sup_{\mu \in \mathcal{M}}\dim_{LY}^P(\rho_\ast \mu) \le \dim_F^P(\rho)\),
was proved in Section \ref{proofLyap}.

\subsection{Proof of \(\dim_F^P(\rho)\le \sup_{\mu \in \mathcal{M}}\dim_{LY}^P(\rho_\ast \mu) \)}
We recall that the limit cone of \(\rho\) is defined as
\[
\mathcal{L}_{\rho} = \lim_{T \to +\infty}\frac{1}{T} (\log S (\rho (\Gamma)),
\]
in the Hausdorff topology. It is a closed convex cone in \(\aaa^+\).
By the \(P\)-Anosov property if \(a \in \mathcal{L}_{\rho} \setminus \lbrace 0\rbrace\) we obtain
\[
\alpha_{i,j}(a) > 0\text{ for all }(i,j) \in S(P).
\]
Assume that \(r < r' < \dim_F^P(\rho)\) so that the series \eqref{falconerseries} diverges at \(r'\). By compactness there exists \(a \in \mathcal{L}_{\rho}\setminus \lbrace 0\rbrace\) such that for all open cones \(\mathcal{C}\) containing \(a\) we have
\[
\sum_{\g \in \G: \log S (\rho(\gamma)) \in \mathcal{C}} \exp \bigl(-F_{r'}^P\circ \log S (\rho(\gamma)) \bigr) = +\infty.
\]
We fix such a choice of \(a\) and notice that \(\psi(a) \ge F_{r'}^P(a)\) by \cite[Lem.\,III.1.3]{Qui02a} and \cite[Lem.\,4.2]{sambarino}. Thus, by Lemma \ref{lyapunovvsfalconer}, \(\psi(a) > F_{r}^P(a)\). We~will prove:
\begin{proposition}\label{goodmeasureslemma}
There exists a sequence \(T_k \to +\infty\) and \(\mu_k \in \M\) such that \(h_{\mu_k} = T_k \psi(a) + o(T_k)\) and \(\lambda(\rho_*\mu_k) = T_k a + o(T_k)\) when \(k \to +\infty\).
\end{proposition}

Since \(\psi(a) > F_r^P(a)\), we have from Lemma \ref{lyapunovvsfalconer} that \(L_{\psi(a)}^P(a) > r\).
Assuming Proposition \ref{goodmeasureslemma}, we obtain
\begin{align*}\dim_{LY}^P(\rho_*\mu_k) &= L_{h_{\mu_k}}^P(\rho_*\mu_k) = L_{T_k\psi(a) + o(T_k)}^P(T_k a + o(T_k))
\\ &= L_{\psi(a) + o(1)}^P(a + o(1)) = L_{\psi(a)}^P(a) + o(1) > r,\end{align*}
for \(k\) large enough. Which concludes the proof of the inequality and of Theorem~\ref{theoremb}.\qed

\subsection{Proof of Proposition \ref{goodmeasureslemma}}

Fix \(\beta:\aaa \to \R\) linear such that \(\beta(a) = F_r^P(a)\), and a decreasing sequence of open cones \(\mathcal{C}_k\), \(k = 1,2,3,\ldots\) whose intersection is \(\R_+ a\).

For a given \(k\) set\vspace*{-3pt}
\[
\psi_k(a) := \limsup_{T \to +\infty}\frac{1}{T}\log \#\bigl\lbrace \gamma \in \Gamma: \log S(\rho(\gamma)) \in \mathcal{C}_k\cap \lbrace \beta \le T \beta(a)\rbrace \bigr\rbrace. 
\]
We have \(\psi (a) = \lim_{k \to \infty } \psi_k (a) \) and\vspace*{-3pt}
\[
\psi_k(a) = \limsup_{T \to +\infty}\frac{1}{T}\log \#\bigl\lbrace \gamma \in \Gamma: \log S(\rho(\gamma)) \in \mathcal{C}_k \cap \lbrace (T-1)\beta(a) \le \beta \le T \beta(a)\rbrace \bigr\rbrace.
\]
Hence, we~may choose \(T_k \to +\infty\) such that the sets\vspace*{-3pt}
\[
A_k = \bigl\lbrace \gamma \in \G: \log S (\rho(\gamma))\in \mathcal{C}_k \cap \lbrace (T_k-1)\beta(a) \le \beta \le T_k \beta(a)\rbrace\bigr\rbrace,
\]
satisfy\vspace*{-3pt}
\[
\psi(a) = \lim_{k \to +\infty}\frac{1}{T_k}\log \# A_k.
\]
Moreover, by the Anosov property, there are constants \(C_1,C_2\) such that for \(\g \in A_k\), \(C_1T_k \le |\g| \le C_2T_k\). So there exist \(\ell_k \to \infty \) as \(k \to \infty \), such that the sets
\[
A'_k = \bigl\lbrace \gamma \in \G: \log S (\rho(\gamma)) \in \mathcal{C}_k \cap \lbrace (T_k-1)\beta(a) \le \beta \le T_k \beta(a)\rbrace,\, |\g| = \ell_k\bigr\rbrace,
\]
still satisfy\vspace*{-3pt}
\[
\psi(a) = \lim_{k \to +\infty}\frac{1}{T_k}\log \# A'_k.
\]

We now fix a semi-group \(S\) of \(\Gamma\), a finite subset \(F\), and maps \(L,R:\Gamma \to F\) given by Proposition \ref{semigrouplemma}. Observe that \(\rho(S)\) is Zariski dense since the Zariski closure of a semi-group is a group(see \cite{goldsheid-margulis}) and since \(S\) generates \(\Gamma\).

Let \(B_k \subset S\) be defined by\vspace*{-3pt}
\[
B_k = \lbrace L(\gamma)\gamma R(\gamma): \gamma \in A'_k\rbrace.
\]
Since \(F\) is finite, there exists \(C > 0\) such that\vspace*{-3pt}
\begin{align*} \# B_k &\ge C^{-1} \# A'_k, & | |\g| - \ell_k | &\le C {\textrm { for }} \g \in B_k \end{align*}
and fixing some norm on \(\aaa\) we have\vspace*{-3pt}
\[
\|\log S(\rho(\gamma)) - \log S(\rho(L(\gamma)\gamma R(\gamma)))\| \le C,
\]
for all \(\gamma \in \Gamma\).

Fix a parameter \(\e> 0 \), we~say a splitting \(\R^d = E_1 \oplus \cdots \oplus E_d\) into one dimensional subspaces is \(\e\)-non-degenerate if
the angle between \(E_i\) and \(\bigoplus_{j \neq i}E_j\) is at least~\(\e\) for all~\(i\). We~say two splittings \((E_1,\ldots, E_d), (E_1',\ldots,E_d')\) are \(\e\)-close if the angle between~\(E_i\) and \(E_i'\) is at most \(\e\) for all \(i\).

We say that an element \(g \in \SL(d,\R)\) is \(\e\)-diagonalizable if it has \(d\) distinct real eigenvalues \(|\lambda_1| > \cdots > |\lambda_d|\) satisfying \(|\lambda_i|/|\lambda_{i+1}| > e^{\e}\) for \(i = 1,\ldots, d-1\), and the corresponding splitting into eigenspaces is \(\e\)-non-degenerate.

Since \(\rho(S)\) is Zariski dense we obtain from \cite[Th.\,6.8]{AMS95}:
\begin{lemma}[Abels-Margulis-Soifer]\label{abslemma}
There exists \(\e_0 > 0\) and a finite set \(S_0 \subset S\) such that for all \(\gamma \in \Gamma\) the set \(\rho(\gamma S_0)\) contains at least one \(\e_0\)-diagonalizable element.
\end{lemma}

For what follows we fix \(S_0\) and \(\e_0 > 0\) given by the previous lemma. We~also fix \(\delta_0 > 0\) given by the following:
\begin{lemma}
There exists \(\delta_0 > 0\) such that if \(B \) is a set of \(\e_0\)-diagonalizable elements of \(\SL(d,\R) \) whose eigenspace splittings are \(\delta_0\)-close, then
\(\gamma_1\cdots \gamma_n\) is \(\e_0/2\)-diagonalizable for all \(n\) and \(\gamma_1,\ldots, \gamma_n \in B\).
\end{lemma}
\begin{proof}
This follows directly from \cite[Lem.\,2.17]{2021breuillard-sert}.
\end{proof}

We now refine the family \(B_k\) so as to obtain images under \(\rho\) which are diagonalizable with nearby splittings.
\begin{lemma}\label{bklemma}
There exists a decreasing positive sequence \(\epsilon_k \downarrow 0\) and a sequence \(\ell '_k \uparrow \infty \) such that for all \(k\) large there exists \(C_k \subset B_k\cdot S_0\) with the following properties:
\begin{enumerate}
\item\label{bklemma1} \(C_k\) freely generates a free semi-group in \(\Gamma\).
\item\label{bklemma2} For all \(\gamma \in C_k\) the element \(\rho(\gamma)\) is \(\e_0\)-diagonalizable.
\item\label{bklemma3} For all \(\gamma_1,\gamma_2 \in C_k\) the eigenspace splittings of \(\rho(\gamma_1)\) and \(\rho(\gamma_2)\) are \(\delta_0\)-close.
\item\label{bklemma4} One has \(\psi(a) = \lim_{k\to +\infty}\frac{1}{T_k}\log \# C_k\).
\item\label{bklemma5} \(|\g| = \ell '_k \) for \(\g \in C_k \).
\item\label{bklemma6} For all \(\gamma \in C_k\) one has \(\|(\log S(\rho(\gamma)) - T_k a\| \le \epsilon_k T_k\).
\end{enumerate}
\end{lemma}
\begin{proof}
The space of \(\e_0\)-non-degenerate splittings is compact and hence we may cover it by some finite number \(N\) of subsets with the property that if two splittings belong to the same subset they are \(\delta_0\)-close.

From Lemma \ref{abslemma} it follows that for all \(\gamma \in B_k\) there exists \(f_\gamma \in S_0\) such that \(\rho(\gamma f_\gamma)\) is \(\e_0\)-diagonalizable. By the pigeonhole principle we may choose a set \(C'_k\) of at least \(\# B_k /N\) of the elements \(\gamma f_{\gamma}\) such that the eigenspace splitting of \(\rho\) applied to any two are \(\delta_0\)-close.
Hence, we~have constructed \(C'_k\) with properties \eqref{bklemma2}, \eqref{bklemma3} and \eqref{bklemma4} above.

Since \(F \cup S_0\) is finite there exists \(C > 0\) such that, for all \(k\) and \(\gamma \in A_k\),
\begin{eqnarray*} | |L(\gamma)\gamma R(\gamma) f_{L(\gamma)\gamma R(\gamma)}| - \ell_k | &\le &C \; {\textrm { and }}\\
\|(\log S(\rho(L(\gamma)\gamma R(\gamma) f_{L(\gamma)\gamma R(\gamma)})) - (\log S(\rho(\gamma))\| &\le &C.\end{eqnarray*}

By the pigeonhole principle again, we~can find \(\ell'_k, | \ell '_k - \ell_k | \le C\), such that \(C''_k : = \g \in C'_k, |\g | = \ell '_k \) satisfies properties \eqref{bklemma2}, \eqref{bklemma3}, \eqref{bklemma4} and \eqref{bklemma5}. Moreover,
if \(\gamma \in A_k\) then by definition we have
\[
\Bigl(\frac{1}{T_k}(\log S(\rho(\gamma))- a\Bigr) \in \mathcal{C}_k \cap \Bigl\lbrace \frac{T_k - 1}{T_k}\beta(a) \le \beta \le \beta(a)\Bigr\rbrace,
\]
hence setting \(\epsilon_k\) to be \(C/T_k\) plus the diameter of the set on the right-hand side, we~have
\[
\|(\log S(\rho(L(\gamma)\gamma R(\gamma) f_\gamma))- T_k a\| \le \epsilon_k T_k,
\]
for all \(\gamma \in A_k\), which establishes property (6) for \(C'_k\) and \(C''_k\).

To establish \eqref{bklemma1} we notice that \(S\) is a quasi-geodesic semi-group. Therefore, by~the Morse lemma, there exists some \(R > 0\) such that if \(\gamma_1,\ldots, \gamma_n, \eta_1,\ldots,\eta_n \in S\) and \(\gamma_1\cdots \gamma_n = \eta_1\cdots \eta_n\), then there is a geodesic ray \(\alpha = \{\alpha_n \}_{ n \ge 0}\), such that for all \(j \in [0, n]\), \begin{equation}\label{tracking}\dist(\gamma_1\cdots \gamma_j, \alpha), \, \dist (\eta_1\cdots \eta_j, \alpha) < R.\end{equation}

We claim that if we assume that the elements of \(C_k \subset C''_k \) are \(6R\)-separated, equation \eqref{tracking} forces \(\g_j = \eta_j \) for all \(j \in [0, n]\).
Hence it suffices to refine \(C''_k\) so that it is a \(6R\)-separated set to obtain that it freely generates a free sub-semi-group of \(\Gamma\). Doing this decreases the number of elements at most by a factor equal to the number of elements in a ball of radius \(6R\) in \(\Gamma\), so property \eqref{bklemma4} still holds.

To prove the claim, we~prove by induction on \(j\) that \(\g_i = \eta_i \) for \(i \le j\). We~have \(\g_0 = \eta_0 = e\). Assume that \(\g_i = \eta_i \) for \(i \le j\). Then
\(\gamma_1\cdots \gamma_j = \eta_1\cdots \eta_j\) and there exist \(N_1, N_2\) and \(N_3\) such that
\[
 \dist (\gamma_1\cdots \gamma_j, \alpha_{N_1}),\;\dist (\gamma_1\cdots \gamma_j\g_{j+1}, \alpha_{N_2}) {\textrm { and }}\dist (\gamma_1\cdots \gamma_j\eta_{j+1}, \alpha_{N_3}) \le R.
\]
Setting \(\beta_t := (\gamma_1\cdots \gamma_j)^{-1} \alpha_{N_t} \) for \(t = 1,2,3\), we have \(\beta_{N_1}, \beta_{N_2}\) and \(\beta_{N_3} \) aligned on the same geodesic and satisfying:
\[
 |\beta_1|,\, \dist (\g_{j+1}, \beta_2),\, \dist (\eta_{j+1}, \beta_3) < R \quand |\g_{j+1}|= |\eta_{j+1}| = \ell'_k.
\]
If \(\ell'_k\) is large enough, \(\beta_1 \) is outside the interval \([\beta_2, \beta_3]\) (see Lemma \ref{boundeddistanelemma}) and
\[
 \dist (\beta_1, \beta_2),\, \dist (\beta_1, \beta_3) \; \in \; [\ell'_k -2R, \ell'_k +2R ].
\]
It follows that \(\dist (\beta_2, \beta_3) \le 4R\) and therefore that \(\dist (\alpha_{j+1}, \eta_{j+1}) < 6R\). Since \(C_k\) is \(6R\)-separated, \(\alpha_{j+1} = \eta_{j+1}\) as claimed.
\end{proof}

\begin{corollary}
If \(\mu_k\) is the uniform probability measure on \(C_k\) then one has
\[
\lim_{k \to +\infty}\frac{1}{T_k}h_{\mu_k} = \psi(a).
\]
\end{corollary}
\begin{proof}
This follows immediately from properties \eqref{bklemma1} and \eqref{bklemma4} of the previous lemma.
\end{proof}

We now use the fact that for \(\e_0\)-diagonalizable elements with close splittings the singular values and eigenvalues are close and they almost multiply under composition.
\begin{lemma}
There exists a decreasing positive sequence \(\epsilon_k' \downarrow 0\) such that for all \(k\) large enough, all \(n\) and \(\gamma_1,\ldots, \gamma_n \in C_k\) one has
\[
\Bigl\|\frac{1}{n}(\log S(\rho(\gamma_1\cdots \gamma_n)) - T_k a\Bigr\| \le \epsilon_k' T_k.
\]
\end{lemma}

\begin{proof}
Let \(L(g)\) denote the vector of absolute values of eigenvalues of a diagonalizable element \(g \in \SL(d,\R)\) in decreasing order.
By \cite[Lem.\,2.16]{2021breuillard-sert} there exists \(C > 0\) such that for all \(\e_0/2\)-diagonalizable \(g \in \SL(d,\R)\) one has
\[
\|\log S(g) - \log L(g)\| \le C.
\]
By \cite[Prop.\,2.17]{2021breuillard-sert} there exists \(C > 0\) independent of \(k\) and \(n\) such that we have
\[
\|\ts (\log L(\rho(\gamma_1\cdots \gamma_n)) - \sum_{i = 1}^n (\log L (\rho(\gamma_i))\| \le Cn.
\]
Combining this with Lemma \ref{bklemma} we obtain
\begin{align*}\|(\log S(\rho(\gamma_1 \cdots \gamma_n)) - n T_k a\| &\le C + \|(\log L (\rho(\gamma_1\cdots \gamma_n))- nT_k a\|
\\ &\ts\le C(n+1) + \|\sum_{i = 1}^n(\log L (\rho(\gamma_i)) - nT_k a\|
\\ &\ts\le C(2n+1) + \|\sum_{i = 1}^n(\log S (\rho(\gamma_i))- nT_k a\|
\\ &\le C(2n+1) + n \epsilon_k T_k.
\end{align*}
This proves the desired result with \(\epsilon_k' = 3C / T_k + \epsilon_k\).
\end{proof}

The following immediate corollary concludes the proof of Proposition \ref{goodmeasureslemma}.

\begin{corollary}
For any sequence of probability measures \(\mu_k\) with \(\mu_k(C_k) = 1\) one has
\[
\lim_{k \to +\infty}\frac{1}{T_k}\lambda(\rho_*\mu_k) = a.
\]
\end{corollary}

\appendix
\section*{Appendix. Quasi-geodesic semi-groups}
\refstepcounter{section}
\subsection{Goal}

If \(\Gamma\) is the free group with generators \(a,b\) then the semi-group \(S\) consisting of reduced words which begin and end with \(a\) has the following nice properties:
\begin{enumerate}
\item \(S\) generates \(\Gamma\) as a group.
\item For every \(n\) and \(\gamma_1,\ldots, \gamma_n \in S\) the geodesic joining the neutral element \(e\) to \(\gamma_1\cdots \gamma_n\) passes through all partial products \(\gamma_1\cdots \gamma_k\) for \(k = 1,\ldots, n\).
\item For every \(\gamma \in \Gamma\) there exist two letter words \(L(\gamma)\) and \(R(\gamma)\) such that \(L(\gamma)\gamma R(\gamma) \in S\).
\end{enumerate}
We may for instance set \(L(\gamma)\) to be \(a\) if the reduced form of \(\gamma\) does not begin with~\(a^{-1}\) and \(ab\) otherwise, and \(R(\gamma)\) similarly to be \(a\) if the last letter of \(\gamma\) is not \(a^{-1}\) and \(ba\) otherwise.

Our goal in this appendix is to give a similar construction for a general hyperbolic group.
For this purpose we fix a group \(\Gamma\) with a finite symmetric generating set \(\cS\). We~let \(|\gamma|\) be the word length of \(\gamma \in \Gamma\) with respect to \(\cS\) and consider the word distance \(\dist(\gamma,\eta) = |\gamma^{-1}\eta|\). We~assume that there exists \(\delta > 0\) (fixed from now on) such that \(\Gamma\) is \(\delta\)-hyperbolic with respect to \(\dist\), and non-elementary (\ie contains two independent loxodromic elements).

Recall that given positive constants \(C,D > 0\) a \((C,D)\)-quasi-geodesic is a sequence \(x: I \cap \Z \to \Gamma\) where \(I\) is an interval, satisfying
\[
C^{-1}|m-n| - D \le \dist(x_m,x_n) \le C|m-n| + D,
\]
for all \(m,n \in I \cap \Z\).

We say a semigroup \(S \subset \Gamma\) is quasi-geodesic if there exist constants \(C,D > 0\) such that for all \(n\) and \(\gamma_1,\ldots, \gamma_n \in S\) the sequence
\[
e, \gamma_1,\gamma_1\gamma_2,\ldots, \gamma_1\cdots \gamma_n,
\]
is visited in left to right order by some \((C,D)\)-quasi-geodesic.

We now state the main result of the appendix.

\begin{proposition}\label{semigrouplemma}
Let \((\Gamma,\cS)\) be a non-elementary hyperbolic group with a finite symmetric generating set \(\cS\). Then there exists a semigroup \(S \subset \Gamma\), a finite set \(F \subset \Gamma\) and mappings \(L,R:\Gamma \to F\) with the following properties:
\begin{enumerate}
\item \(S\) generates \(\Gamma\) as a group,
\item \(S\) is quasi-geodesic, and
\item \(L(\gamma)\gamma R(\gamma) \in S\) for all \(\gamma \in \Gamma\).
\end{enumerate}
\end{proposition}

The arguments in what follows use the local-to-global principle for hyperbolic groups as in \cite[\S 3]{gouezel}.

\subsection{Proof of Proposition \ref{semigrouplemma}}

Recall that the Gromov product (based at the neutral element \(e \in \Gamma\)) is defined by
\[
(\gamma,\eta) = \frac{|\gamma| + |\eta| - |\gamma^{-1}\eta|}{2}.
\]
We assume \(\delta > 0\) is such that
\[
\min\lbrace (\gamma_1,\gamma_2),(\gamma_2,\gamma_3)\rbrace \le (\gamma_1,\gamma_3) + \delta,
\]
for all \(\gamma_1,\gamma_2,\gamma_3 \in \Gamma\). We~further assume that
\[
\dist(e,\alpha) \le (p,q) + \delta,
\]
for all \(p,q \in \Gamma\) and geodesic \(\alpha\) joining them (such a choice of \(\delta\) is possible by \cite[Prop.\,1.22]{bridson-haefliger}).

\begin{lemma}
For each \(C \ge 1\) there exist \(a,b \in \Gamma\) such that setting \(A = \lbrace a,b,a^{-1},b^{-1}\rbrace\) and defining
\[
c_1 = \max_{x,y \in A:\ x \neq y}(x,y)\quand c_2 = \min_{x \in A}|x|,
\]
one has \(c_1 + 2\delta + r < \frac{1}{2}{c_2}\) for some \(r > C(c_1+2\delta+1)\).
\end{lemma}
\begin{proof}
Since \(\Gamma\) is non-elementary we may take \(x,y\) two independent loxodromic elements (\ie the set of boundary fixed points of \(x\) and \(y\) are disjoint).
Letting \(a = x^n\) and \(b = y^n\) for sufficiently large \(n\) establishes the claim since \(c_1(x^n,y^n)\) remains bounded, while \(c_2(x^n,y^n)\) goes to infinity.
\end{proof}

We fix from now on a set \(A\) as in the previous lemma with corresponding constants \(c_1,c_2,r\) satisfying
\begin{equation}\label{requation}
r > 32(c_1+2\delta+1).
\end{equation}

\begin{lemma}\label{borrowedlemma}
The sets defined for \(x \in A\) by \(V(x) = \lbrace \gamma: (\gamma,x) \ge c_1+2\delta + r\rbrace\) are pairwise disjoint.

\begin{itemize}
\item
For all \(x \in A\) and \(\gamma \notin V(x^{-1})\) one has \(x\gamma \in V(x)\).

\item
For all \(x,y \in A\) with \(x \neq y\), if~\(\gamma \in V(x)\) and \(\eta \in V(y)\) then \((\gamma,\eta) \le c_1+2\delta\).

\item
If \(\gamma \in V(x)\) for some \(x \in A\) then \(|\gamma| \ge r\).
\end{itemize}
\end{lemma}
\begin{proof}
By hypothesis \(2\delta + r \in (2\delta,c_2/2 - c_1)\). With this property, the first three statements are claims 1--3 in the proof of \cite[Lem.\,4.1]{random}.

For the last statement observe that
\[
|\gamma| = (\gamma,\gamma) \ge \min\lbrace (\gamma,x),(x,x)\rbrace - \delta \ge \min \lbrace c_1 + 2\delta + r, c_2\rbrace - \delta \ge r.\qedhere
\]
\end{proof}

\begin{lemma}\label{mlemma}
There exists \(m\) such that for all \(\gamma \!\in\! V(a)\) and all \(\eta\) with \(\dist(\gamma,\eta) \!\le\! 2c_2\) one has \(a^m\eta \in V(a)\).
\end{lemma}
\begin{proof}
We first observe that by Lemma \ref{borrowedlemma} we have \((a^{-1},a) \le c_1 + 2\delta\) and therefore a geodesic fixed under multiplication by \(a\) is at distance at most \(c_1+3\delta\) from \(a^m\) for all \(m\).
It follows, letting \(\ell = \lim_{n \to +\infty}\frac{1}{n}|a^n|\) be the translation length of \(a\), that
\[
\ell \ge |a| - 2(c_1+3\delta) \ge 2(c_1+2\delta+r) -2(c_1+3\delta) \ge 2r-2\delta.
\]
Hence we obtain for all \(m \ge 1\) that
\begin{align*}
(a,a^m) &= \frac{|a|+|a^m|-|a^{m-1}|}{2}
\\ &\ge \frac{\ell -2(c_1+3\delta) + m\ell - 2(c_1+3\delta) - (m-1)\ell -2(c_1+3\delta)}{2}
\\ &= \ell - 3(c_1+3\delta)
\ge 2r - 3c_1-11\delta
\ge c_1+3\delta + r.
\end{align*}

In what follows we will need to switch basepoints, so we recall that the Gromov product
\[
(x,y)_z = \frac{\dist(x,z)+\dist(y,z)-\dist(x,y)}{2},
\]
satisfies
\((x,y)_z + (z,y)_x = \dist(x,z)\).
We now calculate using Lemma \ref{borrowedlemma}
\begin{align*}
(a^m,a^{m}\eta) &= (e,\eta)_{a^{-m}}
\\ &= |a^m| - (a^{-m},\eta)_{e}
\\ &\ge |a^m| - (a^{-m},\gamma)_{e} - 2c_2
\ge |a^m| - c_1 - 2\delta - 2c_2
\ge c_1+3\delta + r,
\end{align*}
if we choose \(m\) large enough depending only on \(c_1,c_2\) and \(\delta\).

We conclude the proof using the hyperbolicity property since
\[
(a,a^m\eta) \ge \min\lbrace (a,a^m), (a^m,a^m\eta) \rbrace -\delta \ge c_1 + 2\delta + r,
\]
so \(a^m\eta \in V(a)\) as required.
\end{proof}

We define
\[
T = \lbrace \gamma: \gamma \in V(a)\text{ and }\gamma^{-1} \in V(a^{-1})\rbrace.
\]
With \(m\) given by the previous lemma, for each \(\gamma \in \Gamma\) let
\[
L(\gamma) =\begin{cases}a^{m+1} &\text{if } \gamma \notin V(a^{-1}),\\ a^{m+1}b &\text{otherwise,}\end{cases}\quand
R(\gamma) =\begin{cases}a &\text{if } (L(\gamma)\gamma)^{-1} \notin V(a),\\ ba &\text{otherwise.}\end{cases}
\]

\begin{lemma}
One has \(L(\gamma)\gamma R(\gamma) \in T\) for all \(\gamma \in \Gamma\).
\end{lemma}
\begin{proof}
Since \(\dist(L(\gamma)\gamma, L(\gamma)\gamma R(\gamma)) \le 2c_2\) this follows immediately from Lem\-ma~\ref{mlemma}.
\end{proof}

We now let \(S\) be the semigroup generated by \(T\).

\begin{lemma}
The semigroup \(S\) generates \(\Gamma\) as a group.
\end{lemma}
\begin{proof}
From the definition \(a \in S\) and, using Lemma \ref{borrowedlemma}, \(aba \in S\). Therefore \(L(\gamma),R(\gamma)\) belong to the group generated by \(S\) for all \(\gamma\).
The equation
\[
\gamma = L(\gamma)^{-1}\left(L(\gamma)\gamma R(\gamma) \right)R(\gamma)^{-1}
\]
now implies that \(S\) generates \(\Gamma\) as a group, as required.
\end{proof}

As a first step towards showing that \(S\) is quasi-geodesic we prove the following.

\begin{lemma}\label{boundeddistanelemma}
For any \(n\) and \(\gamma_1,\ldots, \gamma_n \in T\), letting \(\alpha\) be a geodesic joining \(e\) and \(\gamma_1\cdots \gamma_n\), one has\vspace*{-3pt}
\[
|\gamma_1\cdots\gamma_n| \ge |\gamma_1\cdots \gamma_k| + |\gamma_{k+1}\cdots \gamma_n| - r/8,
\]
and \(\dist(\gamma_1\cdots \gamma_k,\alpha) \le r/8\) for all \(k = 0,\ldots, n\).
\end{lemma}
\begin{proof}
Setting \(x_k = \gamma_1\cdots \gamma_k\) for \(k = 0,\ldots, n\), we observe that\vspace*{-3pt}
\[
(x_{k-1},x_{k+1})_{x_k} = (\gamma_{k}^{-1},\gamma_{k+1}) \le c_1 + 2\delta,
\]
and \(\dist(x_k,x_{k+1}) = |\gamma_{k+1}| \ge r\) by Lemma \ref{borrowedlemma}.
Hence the sequence is a \((c_1+\delta,r)\)-chain as in \cite{gouezel}.
We observe that\vspace*{-3pt}
\[
r > 2(c_1 + 2\delta) +4\delta + 1,
\]
and therefore by \cite[Lem.\,3.8]{gouezel} we have
\[
(x_0,x_n)_{x_k} \le c_1 + 4\delta,
\]
for all \(k\).
This implies that\vspace*{-3pt}
\[
\dist(x_k,\alpha) \le c_1 + 5\delta < r/8.
\]
Also since \((x_0,x_n)_{x_k} = (x_k^{-1},x_k^{-1}x_n)_e\) we get by definition
\[
\frac{|\gamma_1\cdots \gamma_k| + |\gamma_{k+1}\cdots \gamma_n| - |\gamma_1\cdots \gamma_n|}{2} \le c_1 + 4\delta \le r/16,
\]
which yields the required lower bound for \(|\gamma_1\cdots \gamma_n|\).
\end{proof}

We now conclude the proof of Proposition \ref{semigrouplemma}.

\begin{lemma}
The semigroup \(S\) is quasi-geodesic.
\end{lemma}
\begin{proof}
Fix \(n\) and \(\gamma_1,\ldots, \gamma_n \in T\), and let \(\alpha\) be a geodesic with \(\alpha_0 = e\) and \(\alpha_{N_n} = \gamma_1\cdots \gamma_n\) with \(N_n = |\gamma_1\cdots \gamma_n|\).
Since the word distance is integer valued, setting \(C = \lfloor r/8 \rfloor\) we have by Lemma \ref{boundeddistanelemma} that there exist \(N_0 = 0, N_1, \ldots,N_n = |\gamma_1\cdots \gamma_n|\) such that\vspace*{-3pt}
\[
\dist(\gamma_1\cdots \gamma_k, \alpha_{N_k}) \le C \le r/8,
\]
for all \(k\).

We also obtain from Lemma \ref{boundeddistanelemma} applied to \(\gamma_1,\ldots, \gamma_{k+1}\) and Lemma \ref{borrowedlemma} that
\begin{align*}
N_{k+1} - N_k &\ge |\gamma_1\cdots \gamma_{k+1}|-r/8 - |\gamma_1\cdots \gamma_k| -r/8
\\ &\ge |\gamma_{k+1}| - 3r/8 \ge 5r/8 \ge 5C,
\end{align*}
and in particular \(N_{k} \le N_{k+1}\) for \(k = 0,\ldots, n-1\).

We now construct a path \(\beta\) which we claim to be quasi-geodesic (with constants independent of \(n\) and \(\gamma_1,\ldots,\gamma_n\)) by concatenating segments of \(\alpha\) with a paths to and from each \(\gamma_1\cdots \gamma_k\) (which will have length at most \(2C\) as seen above).
To be precise, we choose \(\beta:[0,N_n + 2(n-1)C] \cap \Z \to \Gamma\) such that
\begin{enumerate}
\item For each \(k = 0,\ldots,n-1\) one has \(\beta_{m+N_k+2kC} = \alpha_{m+N_k}\) for all \(m \in [0,N_{k+1}-\nobreak N_k] \cap \Z\).
\item \(\beta_{N_k + 2(k-1)C + C} = \gamma_1\cdots \gamma_k\) for \(k = 1,\ldots, n-1\).
\item \(\dist(\beta_m, \beta_{m+1})\le 1\) for all \(m\).
\end{enumerate}
In view of the third property above we have
\(\dist(\beta_l,\beta_m) \le |l-m|\).
For the lower bound, given \(0 \le l < m \le N_n + 2(n-1)C\) we write
\[
\dist(\beta_l,\beta_m) = \sum_{k = 0}^{m-l}\dist(\beta_{l+k},\beta_{l+k+1}) \ge \sum_{k = 0}^{m-l}f(k),
\]
where we set \(f(k) = 1\) if \(k \in [N_i+2iC,N_{i+1}+2iC]\) for some \(i\) and \(f(k) = -1\) otherwise.
Since the sequence \(f(0),\ldots,f(l-m)\) consists of subsequences of runs of at most \(2C\) consecutive times the value \(-1\), with runs of at least \(5C\) times the value~\(1\) in between, we~obtain
\[
\dist(\beta_l,\beta_m) \ge 3(|m-l| - 4C)/7 - 4C,
\]
which establishes that \(S\) is \((7/3, 6C)\)-quasi-geodesic.
\end{proof}


\begin{thebibliography}{99}
\bibitem[AGG+23]{random} A. Azemar, V. Gadre, S. Gouëzel, T. Haettel, P. Lessa \& C. Uyanik – “Random walk speed is a proper function on Teichmüller space”, J. Modern Dyn. 19 (2023), p. 815–832. \url{https://doi.org/10.3934/jmd.2023022}
\bibitem[AMS95]{AMS95} H. Abels, G. A. Margulis \& G. A. Soĭfer – “Semigroups containing proximal linear maps”, Israel J. Math. 91 (1995), no. 1-3, p. 1–30. \url{https://doi.org/10.1007/BF02761637}
\bibitem[Bar10]{barbot} T. Barbot – “Three-dimensional Anosov flag manifolds”, Geom. Topol. 14 (2010), no. 1, p. 153–191. \url{https://doi.org/10.2140/gt.2010.14.153}
\bibitem[BCH10]{BCH10} J. Ban, Y. Cao \& H. Hu – “The dimensions of a non-conformal repeller and an average conformal repeller”, Trans. Amer. Math. Soc. 362 (2010), no. 2, p. 727–751. \url{https://doi.org/10.1090/S0002-9947-09-04922-8}
\bibitem[Ben96]{Ben96} Y. Benoist – “Actions propres sur les espaces homogènes réductifs”, Ann. of Math. (2) 144 (1996), no. 2, p. 315–347. \url{https://doi.org/10.2307/2118594}
\bibitem[BH99]{bridson-haefliger} M. R. Bridson \& A. Haefliger – Metric spaces of non-positive curvature, Grundlehren Math. Wissen., vol. 319, Springer-Verlag, Berlin, 1999. \url{https://doi.org/10.1007/978-3-662-12494-9}
\bibitem[BHR19]{BHR19} B. Bárány, M. Hochman \& A. Rapaport – “Hausdorff dimension of planar self-affine sets and measures”, Invent. Math. 216 (2019), no. 3, p. 601–659. \url{https://doi.org/10.1007/s00222-018-00849-y}
\bibitem[BPS19]{BPS19} J. Bochi, R. Potrie \& A. Sambarino – “Anosov representations and dominated splittings”, J. Eur. Math. Soc. (JEMS) 21 (2019), no. 11, p. 3343–3414. \url{https://doi.org/10.4171/JEMS/905}
\bibitem[BS21]{2021breuillard-sert} E. Breuillard \& Ç. Sert – “The joint spectrum”, J. London Math. Soc. (2) 103 (2021), no. 3, p. 943–990. \url{https://doi.org/10.1112/jlms.12397}
\bibitem[BV19]{BV19} E. Breuillard \& P. P. Varjú – “On the dimension of Bernoulli convolutions”, Ann. Probab. 47 (2019), no. 4, p. 2582–2617. \url{https://doi.org/10.1214/18-AOP1324}
\bibitem[Can]{canary} R. D. Canary – “Anosov representations: informal lecture notes”, https://dept.math. lsa.umich.edu/~canary/lecnotespublic.pdf. 
\bibitem[Dey22]{Dey24} S. Dey – “On Borel Anosov subgroups of $\mathrm{SL}(d, \mathbb{R})$”, 2022, to appear in Geom. Topol., arXiv:2208.02109. \url{https://arxiv.org/abs/2208.02109}
\bibitem[DK22]{DK22} S. Dey \& M. Kapovich – “Patterson-Sullivan theory for Anosov subgroups”, Trans. Amer. Math. Soc. 375 (2022), no. 12, p. 8687–8737. \url{https://doi.org/10.1090/tran/8713}
\bibitem[DO80]{DO80} A. Douady \& J. Oesterlé – “Dimension de Hausdorff des attracteurs”, C. R. Acad. Sci. Paris Sér. A-B 290 (1980), no. 24, p. A1135–A1138. 
\bibitem[Doy88]{Doy88} P. G. Doyle – “On the bass note of a Schottky group”, Acta Math. 160 (1988), no. 3-4, p. 249–284. \url{https://doi.org/10.1007/BF02392277}
\bibitem[DS22]{DS22} L. Dufloux \& V. Suomala – “Projections of Poisson cut-outs in the Heisenberg group and the visual 3-sphere”, Math. Proc. Cambridge Philos. Soc. 172 (2022), no. 1, p. 197–230. \url{https://doi.org/10.1017/S0305004121000177}
\bibitem[Duf17]{Duf17} L. Dufloux – “Hausdorff dimension of limit sets”, Geom. Dedicata 191 (2017), p. 1–35. \url{https://doi.org/10.1007/s10711-017-0240-2}
\bibitem[Fal88]{Fal88} K. J. Falconer – “The Hausdorff dimension of self-affine fractals”, Math. Proc. Cambridge Philos. Soc. 103 (1988), no. 2, p. 339–350. \url{https://doi.org/10.1017/S0305004100064926}
\bibitem[FS22]{FS22} D.-J. Feng \& K. Simon – “Dimension estimates for $C^1$ iterated function systems and repellers. Part II”, Ergodic Theory Dynam. Systems 42 (2022), no. 11, p. 3357–3392. \url{https://doi.org/10.1017/etds.2021.92}
\bibitem[FS23]{FS20} \textemdash{}, “Dimension estimates for $C^1$ iterated function systems and repellers. Part I”, Ergodic Theory Dynam. Systems 43 (2023), no. 8, p. 2673–2706. \url{https://doi.org/10.1017/etds.2022.41}
\bibitem[GdlH90]{GdlH90} É. Ghys \& P. de la Harpe (eds.) – Sur les groupes hyperboliques d’après Mikhael Gromov, Progress in Math., vol. 83, Birkhäuser Boston, Inc., Boston, MA, 1990. \url{https://doi.org/10.1007/978-1-4684-9167-8}
\bibitem[GGKW17]{GGKW17} F. Guéritaud, O. Guichard, F. Kassel \& A. Wienhard – “Anosov representations and proper actions”, Geom. Topol. 21 (2017), no. 1, p. 485–584. \url{https://doi.org/10.2140/gt.2017.21.485}
\bibitem[GM89]{goldsheid-margulis} I. Y. Gol’dshe˘id \& G. A. Margulis – “Lyapunov exponents of a product of random matrices”, Uspehi Mat. Nauk 44 (1989), no. 5(269), p. 13–60. \url{https://doi.org/10.1070/RM1989v044n05ABEH002214}
\bibitem[GMT23]{GMT19} O. Glorieux, D. Monclair \& N. Tholozan – “Hausdorff dimension of limit sets for projective Anosov representations”, J. Éc. polytech. Math. 10 (2023), p. 1157–1193. \url{https://doi.org/10.5802/jep.241}
\bibitem[Gou22]{gouezel} S. Gouëzel – “Exponential bounds for random walks on hyperbolic spaces without moment conditions”, Tunis. J. Math. 4 (2022), no. 4, p. 635–671. \url{https://doi.org/10.2140/tunis.2022.4.635}
\bibitem[Hoc14]{Hoc14} M. Hochman – “On self-similar sets with overlaps and inverse theorems for entropy”, Ann. of Math. (2) 180 (2014), no. 2, p. 773–822. \url{https://doi.org/10.4007/annals.2014.180.2.7}
\bibitem[HR22]{HR19} M. Hochman \& A. Rapaport – “Hausdorff dimension of planar self-affine sets and measures with overlaps”, J. Eur. Math. Soc. (JEMS) 24 (2022), no. 7, p. 2361–2441. \url{https://doi.org/10.4171/jems/1127}
\bibitem[HS17]{HS17} M. Hochman \& B. Solomyak – “On the dimension of Furstenberg measure for $\mathrm{SL}_2(\mathbb{R})$ random matrix products”, Invent. Math. 210 (2017), no. 3, p. 815–875. \url{https://doi.org/10.1007/s00222-017-0740-6}
\bibitem[JLPX24]{JLPX23} Y. Jiao, J. Li, W. Pan \& D. Xu – “On the dimension of limit sets on $\mathbb{P} (\mathbb{R}^3)$ via stationary measures: variational principles and applications”, Internat. Math. Res. Notices (2024), no. 19, p. 13015–13045. \url{https://doi.org/10.1093/imrn/rnae190}
\bibitem[JPS07]{JPS07} T. Jordan, M. Pollicott \& K. Simon – “Hausdorff dimension for randomly perturbed self affine attractors”, Comm. Math. Phys. 270 (2007), no. 2, p. 519–544. \url{https://doi.org/10.1007/s00220-006-0161-7}
\bibitem[Kai00]{Kai00} V. A. Kaimanovich – “The Poisson formula for groups with hyperbolic properties”, Ann. of Math. (2) 152 (2000), no. 3, p. 659–692. \url{https://doi.org/10.2307/2661351}
\bibitem[KLP16]{KLP16} M. Kapovich, B. Leeb \& J. Porti – “Some recent results on Anosov representations”, Transform. Groups 21 (2016), no. 4, p. 1105–1121. \url{https://doi.org/10.1007/s00031-016-9393-6}
\bibitem[KLP17]{KLP17} \textemdash{}, “Anosov subgroups: dynamical and geometric characterizations”, European J. Math. 3 (2017), no. 4, p. 808–898. \url{https://doi.org/10.1007/s40879-017-0192-y}
\bibitem[KLP18a]{KLP18a} \textemdash{}, “Dynamics on flag manifolds: domains of proper discontinuity and cocompactness”, Geom. Topol. 22 (2018), no. 1, p. 157–234. \url{https://doi.org/10.2140/gt.2018.22.157}
\bibitem[KLP18b]{KLP18} \textemdash{}, “A Morse lemma for quasigeodesics in symmetric spaces and Euclidean buildings”, Geom. Topol. 22 (2018), no. 7, p. 3827–3923. \url{https://doi.org/10.2140/gt.2018.22.3827}
\bibitem[KV83]{KV83} V. A. Kaimanovich \& A. M. Vershik – “Random walks on discrete groups: boundary and entropy”, Ann. Probab. 11 (1983), no. 3, p. 457–490. 
\bibitem[KY79]{KY79} J. L. Kaplan \& J. A. Yorke – “Chaotic behavior of multidimensional difference equations”, in Functional differential equations and approximation of fixed points (Bonn, 1978), Lect. Notes in Math., vol. 730, Springer, Berlin, 1979, p. 204–227. \url{https://doi.org/10.1007/BFb0064319}
\bibitem[Käe04]{K04} A. Käenmäki – “On natural invariant measures on generalised iterated function systems”, Ann. Acad. Sci. Fenn. Math. 29 (2004), no. 2, p. 419–458. 
\bibitem[Lab06]{Lab06} F. Labourie – “Anosov flows, surface groups and curves in projective space”, Invent. Math. 165 (2006), no. 1, p. 51–114. \url{https://doi.org/10.1007/s00222-005-0487-3}
\bibitem[Led84]{Led84} F. Ledrappier – “Quelques propriétés des exposants caractéristiques”, in École d’été de probabilités de Saint-Flour, XII–1982, Lect. Notes in Math., vol. 1097, Springer, Berlin, 1984, p. 305–396. \url{https://doi.org/10.1007/BFb0099434}
\bibitem[Lin04]{Lin04} G. Link – “Hausdorff dimension of limit sets of discrete subgroups of higher rank Lie groups”, Geom. Funct. Anal. 14 (2004), no. 2, p. 400–432. \url{https://doi.org/10.1007/s00039-004-0462-y}
\bibitem[LL23]{LL23a} F. Ledrappier \& P. Lessa – “Exact dimension of Furstenberg measures”, Geom. Funct. Anal. 33 (2023), no. 1, p. 245–298. \url{https://doi.org/10.1007/s00039-023-00631-0}
\bibitem[LL24]{LL23} \textemdash{}, “Exact dimension of dynamical stationary measures”, J. Modern Dyn. 20 (2024), p. 679–715. \url{https://doi.org/10.3934/jmd.2024019}
\bibitem[LPX23]{LPX23} J. Li, W. Pan \& D. Xu – “On the dimension of limit sets on $\mathbb{P} (\mathbb{R}^3)$ via stationary measures: the theory and applications”, 2023, arXiv:2311.10265. \url{https://arxiv.org/abs/2311.10265}
\bibitem[MS19]{MS19} I. D. Morris \& P. Shmerkin – “On equality of Hausdorff and affinity dimensions, via selfaffine measures on positive subsystems”, Trans. Amer. Math. Soc. 371 (2019), no. 3, p. 1547–1582. \url{https://doi.org/10.1090/tran/7334}
\bibitem[MS21]{MS21} I. D. Morris \& Ç. Sert – “A strongly irreducible affine iterated function system with two invariant measures of maximal dimension”, Ergodic Theory Dynam. Systems 41 (2021), no. 11, p. 3417–3438. \url{https://doi.org/10.1017/etds.2020.107}
\bibitem[Ose68]{Ose68} V. I. Oseledets – “A multiplicative ergodic theorem. Characteristic Ljapunov, exponents of dynamical systems”, Trudy Moskov. Mat. Obšč. 19 (1968), p. 179–210. 
\bibitem[PSW21]{PSW21} M. B. Pozzetti, A. Sambarino \& A. Wienhard – “Conformality for a robust class of nonconformal attractors”, J. reine angew. Math. 774 (2021), p. 1–51. \url{https://doi.org/10.1515/crelle-2020-0029}
\bibitem[PSW23]{PSW19} \textemdash{}, “Anosov representations with Lipschitz limit set”, Geom. Topol. 27 (2023), no. 8, p. 3303–3360. \url{https://doi.org/10.2140/gt.2023.27.3303}
\bibitem[Qui02]{Qui02a} J.-F. Quint – “Divergence exponentielle des sous-groupes discrets en rang supérieur”, Comment. Math. Helv. 77 (2002), no. 3, p. 563–608. \url{https://doi.org/10.1007/s00014-002-8352-0}
\bibitem[Rap24]{Rap22} A. Rapaport – “On self-affine measures associated to strongly irreducible and proximal systems”, Adv. Math. 449 (2024), article no. 109734 (116 pages). \url{https://doi.org/10.1016/j.aim.2024.109734}
\bibitem[Rue79]{Rue79} D. Ruelle – “Ergodic theory of differentiable dynamical systems”, Publ. Math. Inst. Hautes Études Sci. 50 (1979), p. 27–58. \url{https://doi.org/10.1007/BF02684768}
\bibitem[Sam14]{sambarino} A. Sambarino – “Hyperconvex representations and exponential growth”, Ergodic Theory Dynam. Systems 34 (2014), no. 3, p. 986–1010. \url{https://doi.org/10.1017/etds.2012.170}
\bibitem[Sim15]{Sim15} B. Simon – Harmonic analysis, A comprehensive course in analysis, vol. 3, American Mathematical Society, Providence, RI, 2015. \url{https://doi.org/10.1090/simon/003}
\bibitem[Sul79]{Sul79} D. Sullivan – “The density at infinity of a discrete group of hyperbolic motions”, Publ. Math. Inst. Hautes Études Sci. 50 (1979), p. 171–202. \url{https://doi.org/10.1007/BF02684773}
\bibitem[Var19]{Var19} P. P. Varjú – “On the dimension of Bernoulli convolutions for all transcendental parameters”, Ann. of Math. (2) 189 (2019), no. 3, p. 1001–1011. \url{https://doi.org/10.4007/annals.2019.189.3.9}
\bibitem[Yan19]{yang} W.-y. Yang – “Statistically convex-cocompact actions of groups with contracting elements”, Internat. Math. Res. Notices (2019), no. 23, p. 7259–7323. \url{https://doi.org/10.1093/imrn/rny001}
\bibitem[Zha97]{Zha97} Y. Zhang – “Dynamical upper bounds for Hausdorff dimension of invariant sets”, Ergodic Theory Dynam. Systems 17 (1997), no. 3, p. 739–756. \url{https://doi.org/10.1017/S0143385797085003}
\end{thebibliography}
\end{document}
